\documentclass{article}
\usepackage{natbib} % Required to change bibliography style to APA
\usepackage{amsmath} % Required for some math elements 
\usepackage[utf8]{inputenc}
\usepackage{graphicx}
\usepackage{multicol}
\usepackage[T2A]{fontenc}
\usepackage[russian]{babel}
\usepackage[document]{ragged2e}
\usepackage[left=2.5cm, right=1.5cm, vmargin=2.5cm]{geometry}
\setlength\parindent{0pt} % Removes all indentation from paragraphs

\usepackage[usenames,dvipsnames,svgnames,table]{xcolor}
\usepackage{listings} %% собственно, это и есть пакет listing
\usepackage{caption}
\DeclareCaptionFont{white}{\color{white}} %% это сделает текст заголовка 
\DeclareCaptionFormat{listing}{\colorbox{gray}{\parbox{\textwidth}{#1#2#3}}}
\captionsetup[lstlisting]{format=listing,labelfont=white,textfont=white}
\graphicspath{{images/}}
\usepackage{tikz}
\usetikzlibrary{shapes.geometric, arrows, calc}

%\usepackage{times} % Uncomment to use the Times New Roman font

%----------------------------------------------------------------------------------------
%	DOCUMENT INFORMATION
%----------------------------------------------------------------------------------------

\begin{document}
\lstset{ %
language=[x86masm]Assembler,                 % выбор языка для подсветки (здесь это С)
basicstyle=\small\sffamily, % размер и начертание шрифта для подсветки кода
numbers=left,               % где поставить нумерацию строк (слева\справа)
numberstyle=\tiny,           % размер шрифта для номеров строк
stepnumber=1,                   % размер шага между двумя номерами строк
firstnumber=1,
numberfirstline=true
numbersep=5pt,                % как далеко отстоят номера строк от подсвечиваемого кода
backgroundcolor=\color{white}, % цвет фона подсветки - используем \usepackage{color}
showspaces=false,            % показывать или нет пробелы специальными отступами
showstringspaces=false,      % показывать или нет пробелы в строках
showtabs=false,             % показывать или нет табуляцию в строках
frame=single,              % рисовать рамку вокруг кода
tabsize=2,                 % размер табуляции по умолчанию равен 2 пробелам
captionpos=t,              % позиция заголовка вверху [t] или внизу [b] 
breaklines=true,           % автоматически переносить строки (да\нет)
breakatwhitespace=false, % переносить строки только если есть пробел
escapeinside={\%*}{*)}   % если нужно добавить комментарии в коде
}
\tikzstyle{startstop} = [rectangle, rounded corners, minimum width=2
cm, minimum height=1cm, text width = 3cm,text centered, draw=black, fill=white!30]
\tikzstyle{io} = [trapezium, trapezium left angle=70, trapezium right angle=110,minimum width=3cm, minimum height=1cm, text centered, text width=3cm, draw=black, fill=white!30]
\tikzstyle{process} = [rectangle, minimum width=3cm, minimum height=1cm, text centered, draw=black, fill=white!30]
\tikzstyle{decision} = [diamond, minimum width=3cm, minimum height=1cm, text width=3cm, text centered, draw=black, fill=white!30]
\tikzstyle{arrow} = [thick,->,>=stealth]

\begin{titlepage}


\center % Center everything on the page
 
%----------------------------------------------------------------------------------------
%	HEADING SECTIONS
%----------------------------------------------------------------------------------------

 ФЕДЕРАЛЬНОЕ ГОСУДАРСТВЕННОЕ АВТОНОМНОЕ ОБРАЗОВАТЕЛЬНОЕ УЧЕРЕЖДЕНИЕ ВЫСШЕГО ОБРАЗОВАНИЯ\linebreak  
«Санкт-Петербургский политехнический университет Петра Великого»\\[2cm] % Name of your university/college
\textsc{\Large Институт компьютерных наук и технологий}\\[6.5cm] % Major heading such as course name

%----------------------------------------------------------------------------------------
%	TITLE SECTION
%----------------------------------------------------------------------------------------

{ \huge \bfseries Курсовой проект\\
	\Large \mdseries Программирование на языке Ассемблера\\
	по дисциплине «ЭВМ и переферийные устройства»}\\[6.5cm] % Title of your document

%----------------------------------------------------------------------------------------
%	AUTHOR SECTION
%----------------------------------------------------------------------------------------
\begin{multicols}{2}
\begin{flushright} \large

{Выполнил студент группы: 23504/2:}\\[0.5cm]

{Руководитель\\
профессор, д.т.н.}

\end{flushright}
\begin{flushright}
	
	{Лаппо С.С}\\[0.5cm] % Your nameГ
	
	 
	Молодяков С. А. % Supervisor's Name
	
\end{flushright}
\end{multicols}
% If you don't want a supervisor, uncomment the two lines below and remove the section above
%\Large \emph{Author:}\\
%John \textsc{Smith}\\[3cm] % Your name

%----------------------------------------------------------------------------------------
%	DATE SECTION
%----------------------------------------------------------------------------------------
\flushright{
{\today}\\[0.5cm] % Date, change the \today to a set date if you want to be precise
}
\centering{
Санкт-Петербург\\
2016
}
%----------------------------------------------------------------------------------------
%	LOGO SECTION
%----------------------------------------------------------------------------------------

%\includegraphics{Logo}\\[1cm] % Include a department/university logo - this will require the graphicx package
 
%----------------------------------------------------------------------------------------

\vfill % Fill the rest of the page with whitespace

\end{titlepage}

% Содержание

\tableofcontents
\newpage

% If you wish to include an abstract, uncomment the lines below
% \begin{abstract}
% Abstract text
% \end{abstract}

%----------------------------------------------------------------------------------------
%	SECTION 1
%----------------------------------------------------------------------------------------

\section{Описание работы}

\subsection{Цели}
\begin{enumerate}
\item Дано арифметическое выражение. Разработать программу
проверки правильности по следующим критериям:
\begin{itemize}
	\item допустимые знаки операций $("+" \,,\ "\,-" \,,\ "\,*" \,,\  "\textbackslash")$;
	\item допустимые константы (целые без знака);
	\item допустимые переменные (до пяти латинских букв);
	\item скобки (только круглые, они должны быть парными).
\end{itemize}
Продумать диагностику по разным типам ошибок.
\item Программа "Будильник2". Будильник должен прозвенеть в абсолютное время (11:45). Перехватите прерывание клавиатуры и при повторном наборе (нажатии клавиш). Ваша программа сообщает "Ошибка". При наступлении заданного времени выдайте звуковой сигнал.
 \end{enumerate}

%\subsection{Задачи}
%\begin{enumerate}
%\item Скачивание последней стабильной версии ядра.
%\item Распаковка исходного кода ядра
%\item Конфигурирование ядра, используя прежнюю конфигурацию
%\item Сборка ядра 
%\item Установка нового ядра и проверка успешности установки
%\item Вычисление времени сборки ядра для разного количества потоков(c помощью скрипта: цикл (0..2N+1])
%\item Анализ результатов
%\end{enumerate}

%\subsection{Аппаратная и программная платформы}

%\begin{description}
%\item[Тип устройства]
%Ультрабук
%\item[Производитель]
%Asus
%\item[Модель устройства]
%TP300LD (2014)
%\item[ОЗУ]
%8гб
%\item[Тип накопителя]
%SSD
%\item[Процессор]
%Intel core i7-4510U (8 потоков)
%\item[Операционная система]
%KUbuntu 16.04 LTS
%\end{description} 
 \newpage
%----------------------------------------------------------------------------------------
%	SECTION 2
%----------------------------------------------------------------------------------------

\section{Ход работы}
\subsection{Задача 1}
\flushbottom
\bfseries  Алгоритм \\ \mdseries
Для выполнения работы программа читает вводимые данные и проверяет является ли данный символ константой. Если символ является константой, то она увеличивает счётчик, отвечающий за количество символов в последней константе. Если сивол является одним из знаков операций, если является то проверяет был ли предыдущий таковым. В случае подтверждения предположения - программа выдаёт ошибку и завершается, иначе счётчик отвечающий за количество символов в константе становится равным нулю. Если символ равен точке программа выдаёт ошибку.\\
\begin{tabular}{cc}
	Прерывание &                                Назначение                                \\
	 int 21h   &          В регистре AH 01h - использовалось для чтения символа           \\
	 int 21h   &           В регистре AH 09 - использовалось для вывода строки            \\
	 int 21h   & В регистре AH 4C00h - использовалось для завершения выполнений программы
\end{tabular} \\
\begin{lstlisting}[caption=Код программы]
use16
macro exit_app
{
mov ax,4C00h
int 21h
}
macro print_str str
{
mov ah,9
mov dx,str
int 21h
}

macro read_character
{
mov ah,01h
int 21h
cmp al, 13
je DoneOk
mov cl, bl
mov bl, al
compareToAll
je CheckPrevious
jne UnDone
}

macro compareToAll
{
cmp bl, 42
je CheckPrevious
cmp bl, 43
je CheckPrevious
cmp bl, 47
je CheckPrevious
cmp bl, 45
je CheckPrevious
}

macro compareToAllFP
{
cmp cl, 42
je Fail
cmp cl, 43
je Fail
cmp cl, 47
je Fail
cmp cl, 45
je Fail
cmp cl, 40
je Fail
}

org 100h
mov bx, 1
mov ax, 1
print_str hello

DoneOk:
print_str exit
read_character
exit_app

UnDone:
read_character

CheckPrevious:
compareToAllFP
read_character
Fail:
print_str failMsg
mov ah, 01h
int 21h
exit_app
; data
failMsg db 'Theres fail. Stopping$'
hello db 'Please enter an experssion!$'
exit db 'Looks like everything was ok!$'
\end{lstlisting}
\centering
\begin{tikzpicture}[node distance=2cm]
\node (start) [startstop] {Начало};
\node (in1) [io, below of=start] {Приветственное сообщение};
\draw [arrow] (start) -- (in1);
\node (in2) [io, below of=in1] {Считать очередной символ};
\draw [arrow] (in1) -- (in2);
\node (dec1) [decision, below of=in2, yshift=-2.5cm]{Символ является допустимым, не является вторым подряд знаком, не является знаком препинания, в константе меньше 5 символов};
\draw [arrow] (in2) -- (dec1);
\draw [arrow] (dec1)  -- node[above,pos=0.5] {Да} ($(dec1.east)+ (1,0)$)-- ($(dec1.east) + (1,5.5)$)  -| (in2);
\node (proc1) [io, below of=dec1, yshift=-2.5cm] {Вывести сообщение об ошибке};
\draw [arrow] (dec1) -- node[left,pos=0.5] {Нет} (proc1) ;
\node (stop) [startstop, below of=proc1] {Конец};
\draw [arrow] (proc1) -- (stop);

\end{tikzpicture}\\
Скриншот выполнения:\\
\includegraphics{screen1.png}

\newpage
\subsection{Задача 2}
\flushbottom
\bfseries Алгоритм \\ \mdseries
Для выполнения работы было решено получать системное время, преобразовывать его в строку и сравнивать с заранее заданой строкой времени будильника. При совпадении строк - подать звуковой сигнал и завершить работу программы. Так же для обеспечения работы был написан обработчик прерывания, который при нажатии клавиши увеличивает счётчик, со значением по умолчанию равным 0 на 1, и при достижении определённого значения обработчик вызывает функцию завершенения программы с выводом сообщения об ошибке. \\
\begin{tabular}{cc}
	Прерывание &                                       Назначение                                       \\
	 int 21h   &          В регистре AH 2CH - использовалось для получения системного времени           \\
	 int 21h   &         В регистре AX 3509h - использовалось для получения вектора прерываний          \\
	 int 21h   &         В регистре AH 4CH - использовалось для завершения выполнений программы         \\
	 int 21h   & В регистре AX 2509h - использовалось для установки собственного обработчика прерываний \\
	 int 21h   &          В регистре AH 2, в регистре DL 7 - использовался для генерации звука          \\
	 int 10h   &                       Установка строки в начало (видеорежим DOS)
\end{tabular} \\
\begin{lstlisting}[caption=Код программы]
.MODEL LARGE
code segment
;assume cs:code, ds:code, es:code, ss:code
org 100H

.DATA
PROMPT  DB  'Current System Time is : $',0     
TIME    DB  '00:00:00$',8,0        ; time format hr:min:sec
TIMEEST DB '11:45:00$',8,0
errstr DB 'ERROR$',5,0      

;   Data1   SEGMENT WORD 'DATA'
int1 dw 0h, 0h
x       DB 0
e       DB ?
;   Data1 ends
.CODE 
jmp MAIN

;**************************************************************************;
;-------------------------------  CONVERT  --------------------------------;
;**************************************************************************;

CONVERT PROC 
; this procedure will convert the given binary code into ASCII code
; input : AL=binary code
; output : AX=ASCII code

PUSH DX                       ; PUSH DX onto the STACK 

MOV AH, 0                     ; set AH=0
MOV DL, 10                    ; set DL=10
DIV DL                        ; set AX=AX/DL
OR AX, 3030H                  ; convert the binary code in AX into ASCII

POP DX                        ; POP a value from STACK into DX 

RET                           ; return control to the calling procedure
CONVERT ENDP                   ; end of procedure CONVERT

;**************************************************************************;
;--------------------------------------------------------------------------;
;**************************************************************************;


;**************************************************************************;
;------------------------------  GET_TIME  --------------------------------;
;**************************************************************************;

GET_TIME PROC
; this procedure will get the current system time 
; input : BX=offset address of the string TIME
; output : BX=current time

PUSH AX                       ; PUSH AX onto the STACK
PUSH CX                       ; PUSH CX onto the STACK 

MOV AH, 2CH                   ; get the current system time
INT 21H                       

MOV AL, CH                    ; set AL=CH , CH=hours
CALL CONVERT                  ; call the procedure CONVERT
MOV [BX], AX                  ; set [BX]=hr  , [BX] is pointing to hr
; in the string TIME

MOV AL, CL                    ; set AL=CL , CL=minutes
CALL CONVERT                  ; call the procedure CONVERT
MOV [BX+3], AX                ; set [BX+3]=min  , [BX] is pointing to min
; in the string TIME

MOV AL, DH                    ; set AL=DH , DH=seconds
CALL CONVERT                  ; call the procedure CONVERT
MOV [BX+6], AX                ; set [BX+6]=min  , [BX] is pointing to sec
; in the string TIME

POP CX                        ; POP a value from STACK into CX
POP AX                        ; POP a value from STACK into AX

RET                           ; return control to the calling procedure
GET_TIME ENDP                  ; end of procedure GET_TIME




MAIN PROC
jmp beginn
KeybInt:
cli
push bp
mov bp, sp
push ax
push bx
push cx
push dx
push si

push    ax
in      al,60H             ;read the key
; mov dx, '['
; call print_symbol
; 
;  no, drop through to original driver
pop     ax


in      al,61H             ;get value of keyboard control lines
mov     ah,al              ; save it
or      al,80h             ;set the "enable kbd" bit
out     61H,al             ; and write it out the control port
xchg    ah,al              ;fetch the original control port value
out     61H,al             ; and write it back

mov     al,20H             ;send End-Of-Interrupt signal
out     20H,al             ; to the 8259 Interrupt Controller
;------ other code handles other tests and finally triggers popu
push dxprint_symbol
mov al, x
add al, 1
mov x, al
cmp al, 3
jg exitperr

pop dx
pop si
pop dx
pop cx
pop bx
pop ax
pop bp
sti
iret

beginn: 
mov ax, 3509h
int 21h
mov int1, bx
mov int1+2, es
push cs
pop ds
mov dx, offset KeybInt
mov ax, 2509h
int 21h
pushf
pushf
MOV AX, @DATA                ; initialize DS
MOV DS, AX

LEA BX, TIME                 ; BX=offset address of string TIME

CALL GET_TIME                ; call the procedure GET_TIME
lea si, TIME
lea di, [TIMEEST]
xor dx, dx         ; make sure edx is 0 to start with
LOOPP:
mov al, [si]
mov bl, [di] 
inc si              ; prepare for next char
inc di
cmp al, bl           ; compare two current characters
jne DIFFERENT        ; not equal, get out, you are DONE!
cmp al, 0            ; equal so far, are you at the end?
je SAME              ; got to end of both strings, you're good, get out
jmp LOOPP             ; okay well they agree so far, go to next char
DIFFERENT:
; Do what you need to do for the strings being different
;
jmp DONE
SAME:
; Do what you need to do for the strings being the same
;
;
jmp EXITP
DONE:
;CMP TIME,TIMEEST 

LEA DX, PROMPT               ; DX=offset address of string PROMPT
MOV AH, 09H                  ; print the string PROMPT
INT 21H                      

LEA DX, TIME                 ; DX=offset address of string TIME
MOV AH, 09H                  ; print the string TIME
INT 21H 
mov ax, 0
mov bp, 43690
mov si, 13690
delay2:
dec bp
nop
jnz delay2
dec si
cmp si,0    
jnz delay2

mov ah, 0
mov bh, 0

int 10h     
CALL MAIN
EXITPERR: 

mov ah,9
lea dx,errstr
int 21h     
MOV AH, 4CH                  ; return control to DOS
INT 21H

EXITP:
mov ax,1500
MOV  AH,2     ;??????? ?????? ??????? ?? ?????
MOV  DL,7     ;???????? ??? ASCII 7
INT  21H      ;??????? ?????
MOV AH, 4CH                  ; return control to DOS
INT 21H
MAIN ENDP
END MAIN
\end{lstlisting}
\newpage
\bfseries Cхема алгоритма\\ \mdseries
\centering
\begin{tikzpicture}[node distance=2cm]

\node (start) [startstop] {Начало};
\node (in1) [io, below of=start] {Инициализация данных};
\draw [arrow] (start) -- (in1);
\node (in2) [io, below of=in1] {Получить системное время};
\draw [arrow] (in1) -- (in2);
\node (dec1) [decision, below of=in2, yshift=-1.5cm]{Системное время совпадает с заданным};
\draw [arrow] (in2) -- (dec1);
\draw [arrow] (dec1)  -- node[above,pos=0.5] {Нет} ($(dec1.east)+ (1,0)$)-- ($(dec1.east) + (1,4.5)$)  -| (in2);
\node (proc1) [process, below of=dec1, yshift=-1.5cm] {Проиграть звук};
\draw [arrow] (dec1) -- node[left,pos=0.5] {Да} (proc1) ;
\node (stop) [startstop, below of=proc1] {Конец};
\draw [arrow] (proc1) -- (stop);
\end{tikzpicture}\\
Скриншот выполнения:\\
\includegraphics{screen2.png}




\newpage
\section{Вывод}
\flushleft
В результате выполнения работы были выполнены задания, а так же улучшен навык взаимодействия с системными прерываниями,и создания собственного обработчика прерывания.
%----------------------------------------------------------------------------------------
%	BIBLIOGRAPHY
%----------------------------------------------------------------------------------------

%\bibliographystyle{apalike}

%\bibliography{sample}

%----------------------------------------------------------------------------------------


\end{document}